\documentclass[11pt]{article}
\usepackage{amsmath,hyperref}
\title{A local revoke is not a remote replacement}
\date{}
\begin{document}\maketitle

\paragraph{Scope.}
The preceding native-source fixture allowed atomic control at the
same owner. This develops a different case: the primary queue owner
can revoke, but only a remote authority can perform fallback.
A native request and a native reply can each be delayed or erased.
At most one actual old callback is permitted; the required service
must occur by its allowed contact. This is simulated, not a real
distributed protocol, independently selected mission or field result.

\paragraph{Effect sequence.}
A separate prior report and independent job pay through actual
callbacks. The queued primary can attempt at tick 9 (ticks are
one twentieth of a native time unit), the alternate at tick 11.
An actual revoke request arrives at the original owner at one
of the declared earlier or later times. Its source decision is
irreversible in this fixture; it lacks the alternate actor's
physical authority. A tie can go to either earlier owner according
to the explicit order; once primary has really attempted, this
source cancellation is no longer offered. A genuine revocation
reply is issued only for an actual revoke. The alternative owner
acts only if the reply has actually reached its registered native
callback at or before tick 11. A timely tie gives the reply
precedence. A later received reply is not early authority.

The reply can use an independent physical sender/receiver pair or
either ledger of the two-byte required alternate. The same native
accounting as in the source producer is used, with all tested
effects within one rolling window. The direct sum must leave
enough budget for the requested new work, not just the reply.
No crypto or field message-access claim is made.

\paragraph{Positive and negative cases.}
The optimistic rule revokes on receipt of an early request and
then tries to send its reply. An erased return can strand the
requirement even when the healthy original would have paid.
A late but actually delivered return misses alternate access too.
A received two-byte reply can also use the remaining shared
debit needed by the two-byte alternate. Retaining on insufficient
source evidence avoids a gratuitous loss of a healthy original;
it cannot guarantee rescue if that original will fail. With
qualified reliable timely return and enough actual scoped budget,
early source cancellation does allow the required alternative
to pay. A lost incoming request is not evidence that cancellation
occurred. The test separates the authenticated-law assumption
from the actual native receipt.

In the finite one-request/one-reply catalogue, at the source
decision there can be two indistinguishable continuations:
a healthy primary with lost return, and a failed primary with
a usable return. Revoking in the former misses work; retaining
in the latter misses it. This is not a theorem about arbitrary
communication protocols, replicated logs, pre-staged atomic
alternative authority or future extra observations.

\paragraph{Comparison.}
This controlled boundary resembles the local-information issue in
\href{https://arxiv.org/abs/cs/0006009}{Halpern and Moses}
and the commit question studied by
\href{https://www.microsoft.com/en-us/research/publication/consensus-on-transaction-commit/}
{Gray and Lamport}. No result about those more general protocols
is inferred. The byte-pinned native replay, including actual
erasure and budget effects, agrees with the exact same-information
event reference on 4,608 cases. Three Lean lemmas verify the
limited common-choice obstruction, exclusivity given the guards,
and the charged reply inequality. Loss and field authority
qualification remain explicit. No independent transfer, fault
disposition or programme acceptance follows.
\end{document}
